\answer{ex:13:1}{(a)~$q=100^{1/99}\approx1.048$, passo $\approx4.8\%$; faixa $\approx2.2\cdot10^3$ ($67$~dB). (b)~Passo de $22f_1$, isto é, $220\%$ com $10f_1$.}
\answer{ex:13:2}{$P_{\text{falha}}\approx1.8\cdot10^{-4}$ e $2.8\cdot10^{-16}$: cerca de $300$ falhas por dia com $S=2$ e $5\cdot10^{-10}$ por dia com $S=4$.}
\answer{ex:13:3}{(a)~$H(s)=eF_1T_c/(1+sT_c)^2$, filtro de segunda ordem com frequência de corte $1/(2\pi T_c)\approx3.2$~Hz. (b)~$r\approx22$~disparos/s (pelo primeiro harmônico, $\approx20$).}
\answer{ex:13:4}{(a)~Com $Gd=1/e$ a raiz $s=-1/d$ é dupla, e com nenhum outro $G$ a raiz mais à direita fica mais à esquerda. (b)~$d=30\ms$: $G\approx12$~s$^{-1}$, $\tau=30\ms$; $d=150\ms$: $G\approx2.5$~s$^{-1}$, $\tau=150\ms$: a constante de tempo é igual ao atraso.}
\answer{ex:13:5}{(a)~$E\approx0.427$, $T\approx1.70\,\tau_a=0.68$~s (o modelo com $\tau=20\ms$ dá $0.84$~s). (b)~A equação não contém $I$; o ritmo existe com $b>\beta-1$.}
\answer{ex:13:6}{Com $b=0.8$ não há ritmo; com $b=1.05$, $1.2$, $1.5$, $2$, $4$, $T\approx2.73$, $1.74$, $1.19$, $0.84$, $0.45$~s; com qualquer $I$, $T=0.840$~s: a entrada muda só a amplitude, e não o andamento.}
\answer{ex:13:7}{A substituição $s=u-a$ dá a maior taxa de amortecimento $a+1/d$ com $G=e^{-1-ad}/d$; $a\ge33.3$~s$^{-1}$, $G=e^{-2}/d\approx4.5$~s$^{-1}$; quanto mais forte a co-contração (maior $a$), menor deve ser $G$.}
\answer{ex:13:8}{Problema aberto. Por exemplo, um limiar no acúmulo de fadiga $b[r_i-r_0]_+$ dá $I_{\min}=\beta b r_0/(b-\beta+1)\approx0.53$ com $b=4$, $r_0=0.2$ e período de $1.8$~s com $I=0.6$ até $0.52$~s com $I=3$.}
